% ============================================================
% 第 5 章 问题三：风味指纹与风格条件生成
% ============================================================

\section{问题三：风味指纹与风格条件生成}
\label{sec:q3}

\subsection{问题分析}

问题三包含三个相互关联的子任务：其一，学习以 \texttt{tasting\_category\_fine} 为条件的化合物谱生成模型，使给定风格（如 Tripel、Kriek、Lambic）能生成合理的化合物组合；其二，由生成的化合物谱给出对应的感官画像，即“风味指纹”；其三，完成跨风格迁移，例如把 Pils/Lager 的化合物谱迁移为 Tripel 风格，同时保证 DMS、双乙酰、4VG、反式-2-壬烯醛与 stale hops 等缺陷化合物不恶化。

与问题一、二相比，问题三的困难在于“没有唯一的正确答案”。同一风格的不同啤酒在化合物谱上差异很大，生成模型不能只记住训练样本；同时，风格标签在 \texttt{7.csv} 中分布不均，部分风格样本很少，直接对每个风格单独建模会过拟合。与 GAN/CVAE 类神经生成方法\citep{goodfellow2014generative,mirza2014conditional} 相比，条件高斯 copula 在 $N=250$、$G=22$ 的小样本下参数更少、协方差可收缩，更适合本题；本文因此采用“条件 copula 学风格分布、高斯最优传输做跨风格迁移、style cap 做缺陷约束”的路线，并把评价拆成分布保真、风格保真、多样性、缺陷控制与感官一致五个维度，避免只看生成样本“像不像”。

\subsection{风格标签与风味指纹定义}

以 \texttt{7.csv} 的 \texttt{tasting\_category\_fine} 作为风格标签 $z_i$，通过 \texttt{beer} 列与 \texttt{1.csv} 连接，得到化合物谱—风格配对数据。风格集合记为 $\{1,\dots,G\}$，属于风格 $g$ 的样本集合为 $\Dcal_g$。由于各风格样本量差异较大，本文在训练时采用风格分层划分与协方差收缩，并在结果中单独报告小样本风格的表现。

本文把风格 $g$ 的“风味指纹”定义为其化合物谱均值与感官画像均值的组合：
\begin{equation}
  \mathcal{F}_g=\left(\bar{\vx}_g,\bar{\vy}_g\right),\qquad
  \bar{\vx}_g=\frac{1}{\abs{\Dcal_g}}\sum_{i\in\Dcal_g}\vx_i,\quad
  \bar{\vy}_g=\frac{1}{\abs{\Dcal_g}}\sum_{i\in\Dcal_g}\vy_i.
  \label{eq:q3-fingerprint}
\end{equation}
该定义不是要替代完整的生成分布，而是给生成结果提供一个可比较的“风格中心”：生成样本既要接近 $\bar{\vx}_g$，又要保留风格内部的合理变异，不能退化为均值复制。

\subsection{基线模型：风格均值与全局协方差}

最简单的风格条件生成器是“风格均值 + 高斯噪声”：对每个风格 $g$，估计均值 $\bmu_g$，并用全局协方差 $\bSigma_{\mathrm{pool}}$ 注入随机性：
\begin{equation}
  \tilde{\vx}^{(g)}=\bmu_g+\bSigma_{\mathrm{pool}}^{1/2}\bz,\qquad
  \bz\sim\Ncal(\bzero,\bI).
  \label{eq:q3-mean-noise}
\end{equation}
第二个基线是“全局协方差 + 风格均值平移”：保留全局相关结构，只把均值移到目标风格。它们可复现、训练快，能给出风格中心与基本变异，但无法刻画风格内协方差的差异，也容易在长尾/零膨胀化合物上生成不合理的极端值；下面的条件高斯 copula 针对这一点改进。

\subsection{改进模型：条件高斯 copula}

化合物浓度具有强零膨胀与长尾，直接用高斯模型会生成负值或极端值。本文先把每个化合物 $j$ 的经验分布 $F_j$ 映射到标准正态分数
\begin{equation}
  z_{i,j}=\Phi^{-1}\left(F_j(x_{i,j})\right),
  \label{eq:q3-normal-score}
\end{equation}
其中 $F_j$ 由训练集的经验分布（或分位数变换）估计，$\Phi$ 为标准正态累积分布函数。对风格 $g$，估计均值 $\bmu_g$ 与协方差 $\bSigma_g$；当 $\abs{\Dcal_g}$ 较小时用收缩估计
\begin{equation}
  \bSigma_g^{(\lambda)}=(1-\lambda)\hat{\bSigma}_g+\lambda\bSigma_{\mathrm{pool}},
  \label{eq:q3-shrinkage}
\end{equation}
其中 $\bSigma_{\mathrm{pool}}$ 为所有风格合并后的协方差，$\lambda\in[0,1]$ 由交叉验证或 Ledoit--Wolf 规则选择。生成时
\begin{equation}
  \bz\sim\Ncal(\bmu_g,\bSigma_g^{(\lambda)}),\qquad
  \tilde{x}_j=F_j^{-1}\left(\Phi(z_j)\right),
  \label{eq:q3-copula}
\end{equation}
并对 $\tilde{x}_j$ 做非负截断。式~\eqref{eq:q3-copula} 是一个条件高斯 copula：它保留每个化合物的边际分布与风格内的相关结构，同时把风格信息压缩为均值与协方差。与风格均值/全局协方差基线相比，收缩协方差在小样本风格上更稳定；与条件变分自编码器（CVAE）相比，它没有神经网络的高方差，更适合 $N=250$、$G=22$ 的数据规模。本文把 CVAE 作为后续扩展，而不是当前主模型\citep{kingma2014auto,sohn2015learning}。

\subsection{最终模型：缺陷约束生成与最优传输迁移}

\paragraph{缺陷约束。}记缺陷化合物集合为 $\Jcal$（至少包含 DMS、双乙酰、4VG、反式-2-壬烯醛；stale hops 用代理变量），其阈值为 $\tau_j$。定义风格感知上限
\begin{equation}
  \bar{c}_{g,j}=\max\left(\tau_j,\ q_{g,j}\right),
  \label{eq:q3-cap}
\end{equation}
其中 $q_{g,j}$ 为风格 $g$ 中化合物 $j$ 的 75\% 分位数。生成后对缺陷列做投影
\begin{equation}
  \tilde{x}_j\leftarrow\min\left(\tilde{x}_j,\bar{c}_{g,j}\right),\qquad j\in\Jcal.
  \label{eq:q3-projection}
\end{equation}
这样既不把缺陷一律压到零（避免破坏风格基线），也不允许生成超出风格正常波动的极端值；同时报告严格阈值 $\tau_j$ 下的超阈值率，区分“风格允许”与“绝对超标”。

其中 4VG 的风格上限取固定安全值 $0.10$~mg/L 而非 $\max(\tau_j,q_{g,j})$，因为小麦风格 4VG 的 $p_{75}$ 显著高于该安全值；其余三个缺陷按式~\eqref{eq:q3-cap} 取风格上限。stale hops 不是单一化合物，本文用反式-2-壬烯醛、DMS、\texttt{ethyl 2-methylbutyrate}、\texttt{ethyl isovalerate} 与 4VG 五列标准化均值构造代理，并在迁移中对每一列施加“不高于源样本”的约束。各缺陷的阈值与依据见表~\ref{tab:q3-thresholds}。

\begin{table}[htbp]
  \centering
  \caption{缺陷化合物的严格阈值与取值依据}
  \label{tab:q3-thresholds}
  \footnotesize
  \setlength{\tabcolsep}{3pt}
  \begin{tabular}{@{}l l l l l@{}}
    \toprule
    缺陷 & 对应列 & 严格阈值 $\tau_j$ & 单位 & 依据 \\
    \midrule
    DMS & \texttt{DMS} & 17.25 & \textmu g/L & 附件与文献综合 \\
    Diacetyl & \texttt{X2.3.butanedione} & 0.09 & mg/L & 附件与文献综合 \\
    4VG & \texttt{X4VG} & 0.10 & mg/L & 非小麦安全值（均值 0.206） \\
    trans-2-nonenal & \texttt{X.E..2.nonenal} & 0.10 & \textmu g/L & 外部文献 \\
    stale hops & 五列复合代理 & — & — & 见正文说明 \\
    \bottomrule
  \end{tabular}
\end{table}

\paragraph{跨风格迁移。}设源风格为 $g$、目标风格为 $g'$。在正态分数空间，用高斯最优传输的闭式映射
\begin{equation}
  T_{g\to g'}(\bz)=\bmu_{g'}+\bSigma_{g'}^{1/2}\bSigma_g^{-1/2}\left(\bz-\bmu_g\right),
  \label{eq:q3-ot}
\end{equation}
将源样本映射到目标风格，再逆变换回浓度空间并做式~\eqref{eq:q3-projection} 的缺陷投影。式~\eqref{eq:q3-ot} 是均值和协方差匹配意义下的最优传输映射；它比“只加风格均值差”多了协方差校正，能同时改变化合物的相关结构。对 Pils/Lager $\to$ Tripel，本文比较三种迁移：raw（只做 OT，不加缺陷约束）、strict（投影到附件严格阈值）、style（投影到式~\eqref{eq:q3-cap} 的风格上限），并报告三者的缺陷超标率与感官画像。

\paragraph{目标边际校准。}仿射 OT 只匹配均值与协方差，不保证各化合物的边际形状（零膨胀的痕量化合物尤甚）。因此迁移后对每个化合物列再做一次秩映射，把生成样本的分位数对齐到目标风格的经验分布：
\begin{equation}
  \tilde{x}_{i,j}\leftarrow Q_{g',j}\!\left(\frac{\operatorname{rank}(\tilde{x}_{i,j})-1/2}{n}\right),
  \label{eq:q3-rank}
\end{equation}
其中 $Q_{g',j}$ 为目标风格 $g'$ 化合物 $j$ 的经验分位数函数。该步之后每一列的边际与目标风格严格一致，因此后文的 MMD 与覆盖度衡量的是\textbf{化合物之间的相关结构}，而非整体分布形状；引用 $0.038$ 这一数值时必须带上此口径。

\begin{algorithm}[htbp]
  \small
  \caption{缺陷约束的最优传输风格迁移}
  \label{alg:transfer}
  \Input{源样本 $\vx_i$、源风格 $g$、目标风格 $g'$、边际分布 $F_j$、阈值 $\tau_j$、风格分位数 $q_{g',j}$}
  \Output{迁移后的化合物谱 $\tilde{\vx}_{\mathrm{tr}}$}
  计算源/目标风格的正态分数均值与协方差 $\bmu_g,\bSigma_g,\bmu_{g'},\bSigma_{g'}$\;
  对源样本做正态分数变换 $\bz_i=\Phi^{-1}(F(\vx_i))$\;
  用式~\eqref{eq:q3-ot} 映射到目标风格，再逆变换回浓度空间\;
  \For{每个缺陷化合物 $j\in\Jcal$}{
    投影 $\tilde{x}_j\leftarrow\min(\tilde{x}_j,\max(\tau_j,q_{g',j}))$\;
  }
  输出 $\tilde{\vx}_{\mathrm{tr}}$\;
\end{algorithm}

\subsection{实验设置与评价指标}

本文按风格分层划分训练/验证/测试集，保证每个风格在测试集中至少有一个样本；小样本风格采用收缩估计，并在结果中单独标注其不确定性。生成模型的对照包括：风格均值 + 高斯噪声（mean\_noise）、全局协方差 + 风格均值平移（global\_cov）、条件高斯 copula、copula + 缺陷约束，以及 OT 迁移。所有模型使用相同的训练/测试划分与随机种子。Tripel 28、Kriek 10、Lambic 11 个样本；每风格每模型生成 300 条；输入为 70 个校准化合物，做 log1p 与逐列标准化；协方差用 Ledoit--Wolf 收缩；OT 后做逐特征目标分位数校准；缺陷投影用 $\max(\tau_j,q_{g,j})$；感官代理用 11 成分 PLS。

评价从五个维度展开：

\begin{enumerate}
  \item \textbf{风格保真}：用真实数据训练风格分类器，在生成样本上的分类准确率；同时计算生成样本到目标风格中心的马氏距离
  \begin{equation}
    D_g(\tilde{\vx})=\sqrt{\left(\tilde{\vx}-\bmu_g\right)^\top
    \bSigma_g^{-1}\left(\tilde{\vx}-\bmu_g\right)}.
    \label{eq:q3-mahalanobis}
  \end{equation}
  \item \textbf{多样性}：生成样本两两之间的平均距离，以及各化合物维度的方差；防止模型只输出风格均值。另有覆盖度与精确度两个派生指标：以真实样本留一最近邻距离的中位数为阈值，覆盖度定义为“至少被一个生成样本落入该阈值的真实样本比例”；以真实留一最近邻距离的 95\% 分位为阈值，精确度定义为“到真实集最近邻距离不超过该阈值的生成样本比例”。
  \item \textbf{缺陷控制}：缺陷化合物超阈值比例
  \begin{equation}
    R_{\mathrm{def}}=\frac{1}{\abs{\Jcal}}
    \sum_{j\in\Jcal}\mathbf{1}\left\{\tilde{x}_j>\tau_j\right\}.
    \label{eq:q3-defect-rate}
  \end{equation}
  \item \textbf{感官一致性}：用问题一最终模型预测生成谱的感官画像 $\hat{\vy}(\tilde{\vx})$，与目标风格画像 $\bar{\vy}_g$ 比较
  \begin{equation}
    \mathrm{MAE}_{\mathrm{sen}}=\frac{1}{S}\sum_{s=1}^{S}
    \abs{\hat{y}_s(\tilde{\vx})-\bar{y}_{g,s}}.
    \label{eq:q3-sensory}
  \end{equation}
\end{enumerate}

此外，用 RBF 核的最大均值差异（MMD）衡量生成分布与真实风格分布的整体差异：
\begin{equation}
  \mathrm{MMD}^2=
  \frac{1}{n^2}\sum_{i,i'}k(\vx_i,\vx_{i'})
  +\frac{1}{m^2}\sum_{j,j'}k(\tilde{\vx}_j,\tilde{\vx}_{j'})
  -\frac{2}{nm}\sum_{i,j}k(\vx_i,\tilde{\vx}_j),
  \label{eq:q3-mmd}
\end{equation}
其中 $k(\cdot,\cdot)$ 为 RBF 核。五个维度必须同时报告：只提高风格分类准确率可能意味着模式崩塌，只提高多样性可能生成不合理化合物，只压低缺陷比例可能把风味一起压没。本文用“分布—保真—多样—缺陷—感官”五象限表呈现结果，避免单一指标误导。

\subsection{结果与分析}

表~\ref{tab:q3-results} 给出三风格平均的五维评价。OT+约束在分布保真（MMD 0.038）、覆盖度（0.863）与 top-1 命中率（0.981）上均最好；条件 copula 的命中率 0.974 接近，但 MMD（0.111）与感官 RMSE（0.132）更差，说明它更接近风格中心、却牺牲了分布与感官一致性；均值噪声的命中率 0.851、覆盖度 0.300，存在模式崩塌。缺陷列显示，OT+约束把 T2N/4VG 的超严格阈值率降到 0，但 DMS/Diacetyl 仍分别有 0.59/0.81 的超严格阈值率，这是因为 Tripel/Kriek/Lambic 的风格 p75 本身高于附件严格阈值；相对 style cap 的超标率为 0（四缺陷均 0.000），详见迁移段落。因此“风格允许”与“绝对安全”必须分开报告；表中的超严格阈值率按 DMS/Diacetyl/T2N 三项报告。

\figauto{q3_mmd_classifier.pdf}{生成模型的分布距离与风格分类命中率}{fig:q3mmd}

\figauto[1.0\textwidth]{q3_transfer_defects.pdf}{Pils/Lager $\to$ Tripel 迁移的缺陷控制}{fig:q3transfer}

迁移结果方面，源风格 Pils/Lager 的 Tripel 分类命中率为 0.00；OT 迁移后 raw/strict/style 三种变体的 Tripel 命中率均为 0.91（随机森林），感官画像与 Tripel 的 Pearson 相关分别为 0.974/0.856/0.972，说明风格方向迁移成功。raw 迁移把 DMS/Diacetyl/T2N 从 15.8 \textmu g/L、0.117 mg/L、0.004 \textmu g/L 提高到 41.8、0.418、0.115，11/11 个样本均“恶化”；strict 投影把三者压到 13.7、0.047、0.004，超严格阈值率为 0，且相对源样本的恶化比例为 0；style 投影保留 Tripel 的风格基线（39.8、0.398、0.099），相对 Tripel p75 的超标率为 0，但相对 Pils 的升高是风格本身的特征，不应误判为缺陷恶化。关键化合物均值从 Pils 的乙醇 5.19\%、甘油 1.64 g/L、乙酸乙酯 10.8 mg/L、乙酸异戊酯 1.15 mg/L、异戊醇 27.9 mg/L，迁移到 Tripel 的 8.79\%、3.11 g/L、40.3 mg/L、4.62 mg/L、54.6 mg/L；linalool 均值从 0.062 升到 1.14 mg/L，高于 Tripel 中位数 0.372 mg/L 但低于其均值 2.97 mg/L，说明 OT 在长尾化合物上可能过冲，后续应加分位数约束。整体上，“绝对阈值”与“风格基线”两种口径给出的结论不同：前者更保守，后者更符合风格真实性；本文同时报告两者，不把“风格允许”误写成“绝对安全”。

\begin{table}[htbp]
  \centering
  \caption{问题三生成模型的五维评价（三风格平均）}
  \label{tab:q3-results}
  \footnotesize
  \setlength{\tabcolsep}{3pt}
  \begin{tabularx}{\textwidth}{@{}l X c c c c@{}}
    \toprule
    评价维度 & 指标 & 均值噪声 & 全局协方差 & 条件 copula & OT+约束 \\
    \midrule
    分布保真 & $\mathrm{MMD}^2$（越低越好） & 0.063 & 0.066 & 0.111 & 0.038 \\
    风格保真 & 随机森林 top-1 命中率 & 0.851 & 0.667 & 0.974 & 0.981 \\
    多样性 & 覆盖度（越高越好） & 0.300 & 0.233 & 0.430 & 0.863 \\
    缺陷控制 & 超严格阈值率 & 0.61/0.67/0.40 & 0.57/0.60/0.43 & 0.56/0.73/0.07 & 0.59/0.81/0.00 \\
    感官一致 & 与目标画像的 $\mathrm{RMSE}$ & 0.090 & 0.095 & 0.132 & 0.090 \\
    \bottomrule
  \end{tabularx}
\end{table}

\figauto[1.0\textwidth]{q3_radar.pdf}{各风格风味指纹的感官画像（雷达图示例）}{fig:q3radar}

图~\ref{fig:q3radar} 用雷达图对比各风格的感官画像。雷达图的每个轴对应一个感官维度，取值来自问题一模型的预测或真实样本均值；它用于检查生成样本是否在多个感官轴上同时接近目标风格，而不是只在单一轴上“像”。

案例部分重点展示三类结果：其一，给定 Tripel 生成的代表性化合物谱与感官画像，检查其是否呈现预期的果香/醇厚特征；其二，Kriek 与 Lambic 的对比，检查模型是否区分了风格间的关键差异，而不是只学到“都是酸啤”；其三，Pils/Lager 到 Tripel 的迁移，检查缺陷化合物是否仍在阈值内、感官画像是否向 Tripel 移动。案例结果为：Tripel/Kriek/Lambic 的生成画像与真实风格中位数的 Pearson 分别为 0.968/0.957/0.990，RMSE 0.051/0.147/0.071；Kriek 的 RMSE 较大，反映其真实感官分布更分散。Pils $\to$ Tripel 的 strict 迁移在保持 Tripel 命中率 0.91 的同时把四类缺陷超严格阈值率压到 0，style 迁移保留 0.972 的感官相关但保留风格基线缺陷；两者分别对应“安全优先”和“风格优先”两种使用场景。

\subsection{讨论}

问题三的评价没有唯一真值，本文用五维指标与案例结合，但仍有三点局限。第一，风格分类器本身是在真实数据上训练的，若真实数据存在标签噪声，生成样本的“保真”分数会继承这种噪声；因此本文同时报告马氏距离与感官一致性，避免只依赖分类器。第二，缺陷阈值来自 \texttt{8.xlsx} 与题目给定口径，不同风格、不同市场的阈值可能不同；本文把阈值作为可配置参数，并在敏感性分析中改变阈值，检查结论是否稳健。第三，OT 映射假设正态分数空间中的风格差异可由均值/协方差线性刻画；若风格流形高度非线性，线性 OT 可能只迁移了部分特征。后续可改用核 OT 或条件流模型，但必须保留缺陷约束与可解释性检查。

另有两处口径必须在结论前交代。其一，\texttt{X.E..2.nonenal} 在附件中标注为 mg/L，但其数值范围（Pils/Lager 中位数约 0.002、Tripel 约 0.011）与文献常用的 \textmu g/L 阈值量级一致，本文按 \textmu g/L 解释；若按标注单位计算，全部样本都远低于检测限，分类将失去意义。其二，stale hops 代理中的 \texttt{ethyl 2-methylbutyrate} 与 \texttt{ethyl isovalerate} 本身属于果香酯类，把它们纳入“不恶化”约束会同时抑制 Tripel 应有的酯类上升，因此 strict 迁移对这两列是\textbf{保守偏置}，不能读作对 stale hops 的独立实测。

另一个需要如实说明的问题是数据规模。部分风格样本很少，生成模型在该风格上的多样性指标可能被高估或低估；本文不把少数风格的结果外推到所有商业啤酒，而是把它们作为“数据条件下的可行方案”，并建议后续补充同风格样本后再复评。
